We introduce \emph{Morphogenic Intelligence} (MI), a neural architecture in which compute units --- called \emph{Morphons} --- grow, differentiate, fuse, and undergo apoptosis at runtime. Unlike Transformers, CNNs, RNNs, and SSMs, where topology is fixed at design time and only weights are learned, MI treats network architecture itself as a primary learned artifact. Each Morphon carries a developmental program inspired by biological morphogenesis: chemical gradients in a hyperbolic embedding space guide connectivity, neuromodulatory signals gate plasticity, and a dual-clock scheduler separates fast inference from slow structural remodeling. We implement the full MI engine in Rust and evaluate it on CartPole control and MNIST classification. CartPole-v1 is solved (avg=195.2 on the v0.5.0 codebase) using only local three-factor learning and developmental morphogenesis. On MNIST, the system reaches \textbf{52\%} intact accuracy on the standard profile --- a 22 percentage point improvement over the v3.0.0 baseline, achieved by replacing rate-based inhibitory STDP with the spike-timing rule of \citet{vogels2011inhibitory}, which breaks hub-morphon dominance (from 1 labeled class to 12). An architectural optimization --- replacing array-of-structs Morphon storage with cache-resident hot arrays (\emph{Pulse Kernel Lite}) --- delivers a \textbf{4.3$\times$} wall-clock speedup on the MNIST workload with no accuracy loss. We document and fix four characteristic failure modes (modulatory explosion, motor silencing, LTD vicious cycle, premature regulatory throttling), each of which has a direct biological analogue. The premature throttling bug --- where the homeostatic regulator declares a system ``mature'' at 11\% accuracy because dense supervised reward inflates its slow EMA --- is novel and broadly applicable to any biologically-inspired regulator that uses raw reward as a maturity signal. A post-training diagnostic (v4.4.0) places MORPHON in the Liquid State Machine (LSM) regime: offline classifiers on frozen morphon activation vectors achieve only 10.5\% --- random chance --- while a sensory-only readout ablation shows that associative morphons contribute \textbf{35.5pp} of the 48.5\% online accuracy. The reservoir computes genuine discriminative information, but the features are co-adapted to the learned readout weights and are not offline-extractable without the dynamical context of the running system. The accuracy ceiling at $\sim$50\% is accordingly a co-adaptation boundary, not a capacity or tuning limit. The primary contributions of this work are not the accuracy numbers --- 52\% on MNIST is well below supervised baselines --- but the failure-mode taxonomy and the architectural characterizations it produces. The premature Mature bug, the LSM representational regime, and the principle that damage-induced recovery in a developmental system signals a regulatory failure rather than a useful property are findings that transfer beyond MI to any biologically-inspired learning architecture. While MI does not yet match static architectures on classification accuracy, it demonstrates that developmental dynamics are a principled and underexplored axis of neural architecture design, and that biology offers not merely metaphors but actionable engineering fixes.
